\section{An agenda for advantage-aware evaluation}
\label{sec:future}

If the gap is that no benchmark is built to ask whether prediction helps acting, the remedy is a
small, measurable protocol that any benchmark could adopt. Table~\ref{tab:future_metrics} sets out
four quantities, each with a definition, a concrete testbed, and the axis it stresses, and
Figure~\ref{fig:future_protocol} sketches how they compose. The figure shows target shapes, not
measured results: it is a protocol, not a leaderboard.

The first quantity is an \emph{advantage-of-prediction curve}: the closed-loop success gain of a
world-model policy over a matched direct-VLA baseline, reported per capability as a curve rather than
a single endpoint, instantiated over LIBERO~\cite{libero} and CALVIN~\cite{calvin} run under one
harness. The second is \emph{counterfactual accuracy}: task success on episodes that require
reasoning over unobserved or hypothetical dynamics, built from WorldPrediction~\cite{worldprediction}
together with occlusion and memory splits in the style of LIBERO-Mem~\cite{liberomem}; this targets
the capability the corpus most neglects. The third is the \emph{prediction-to-action fidelity gap}:
the difference between a world model's open-loop generation quality and its closed-loop task success,
measured on bridge protocols such as WorldSimBench~\cite{worldsimbench} and
World-in-World~\cite{worldinworld}, which quantifies directly the ``visual quality is not task
success'' warning of Section~\ref{sec:gap}. The fourth is \emph{per-capability contrast coverage}:
the fraction of capabilities for which a benchmark actually runs a VLA-versus-world-model head-to-head,
computed by auditing the landscape (Table~\ref{tab:landscape}) against a declared contrast schema.
Each is reported as a distribution or a curve, never as one number, so that ``prediction helps'' can
be replaced by ``prediction helps \emph{here}, by \emph{this much}''.

% Future-agenda metrics: metric / definition / testbed / capability-family stressed.
\begin{table*}[t]
\centering
\footnotesize
\setlength{\tabcolsep}{4pt}
\renewcommand{\arraystretch}{1.35}
\resizebox{\textwidth}{!}{%
\begin{tabular}{@{}p{0.19\textwidth} p{0.30\textwidth} p{0.27\textwidth} p{0.15\textwidth}@{}}
\toprule
\textbf{Metric} & \textbf{Definition} & \textbf{Instantiation (testbed)} & \textbf{Axis stressed} \\
\midrule
Advantage-of-prediction curve &
closed-loop success \emph{gain} of a predictive policy over a matched direct-VLA baseline,
reported per capability as a curve, not a single endpoint &
a capability-sliced protocol over LIBERO~\cite{libero} and CALVIN~\cite{calvin} that runs a
world-model policy and a VLA policy under one harness &
$\gamma$ family, $\delta$ metric \\
Counterfactual accuracy &
task success on episodes that require inference over unobserved or hypothetical dynamics
(occlusion, object permanence, ``what if'') &
WorldPrediction~\cite{worldprediction} plus occlusion / memory splits in the style of
LIBERO-Mem~\cite{liberomem} &
$\beta$ capability \\
Prediction-to-action fidelity gap &
the difference between a world model's open-loop generation quality and its closed-loop task
success (``visual quality is not task success'') &
bridge protocols such as WorldSimBench~\cite{worldsimbench} and
World-in-World~\cite{worldinworld} &
$\alpha$ mode \\
Per-capability contrast coverage &
the fraction of capabilities for which a benchmark actually runs a VLA-vs-world-model
head-to-head under a single protocol &
audits of the landscape (Table~\ref{tab:landscape}) against a declared contrast schema &
$\gamma$ family, $\varepsilon$ catalog \\
\bottomrule
\end{tabular}}
\caption{\textbf{An actionable protocol for the future agenda}: four measurable quantities, their
definitions, and concrete testbeds. The first anchors the advantage-of-prediction question
($\delta$) on LIBERO / CALVIN-style suites; the second targets the near-unmeasured counterfactual
capability; the third quantifies the prediction-to-action gap the bridges expose; the fourth audits
per-capability contrast coverage across the landscape. Each is reported as a distribution or curve,
never a single headline number.}
\label{tab:future_metrics}
\end{table*}

% Future evaluation-protocol figure. SCHEMATIC / ILLUSTRATIVE ONLY: the panels show the axes and the
% target shape of each measurement, NOT measured results. Pairs with Table~\ref{tab:future_metrics}.
\begin{figure}[t]
\centering
\resizebox{\textwidth}{!}{%
\begin{tabular}{@{}c@{\hspace{9mm}}c@{}}
% ---- (a) advantage of prediction ----
\begin{tikzpicture}[baseline]
  \draw[->,gray] (0,0)--(3.8,0) node[right,font=\scriptsize\bfseries,text=black]{capability};
  \draw[->,gray] (0,0)--(0,2.4);
  \node[rotate=90,font=\scriptsize\bfseries,anchor=south,text=gray!10!black] at (-0.2,1.25){closed-loop success};
  \fill[secE!12] (0,1.25) .. controls (1.5,1.85) and (2.4,1.7) .. (3.5,1.65)
                 -- (3.5,0.8) .. controls (2.4,0.95) and (1.5,1.0) .. (0,0.9) -- cycle;
  \draw[gray,line width=1pt,dashed] (0,0.9) .. controls (1.5,1.0) and (2.4,0.95) .. (3.5,0.8);
  \draw[secE!70!black,line width=1.3pt] (0,1.25) .. controls (1.5,1.85) and (2.4,1.7) .. (3.5,1.65);
  \node[secE!75!black,font=\scriptsize\bfseries,anchor=west] at (1.0,1.95){world-model policy};
  \node[gray!10!black,font=\scriptsize\bfseries,anchor=west] at (1.1,0.62){direct VLA baseline};
  \node[font=\footnotesize\bfseries] at (1.9,2.75){(a) Advantage of prediction ($\delta$)};
  \node[font=\scriptsize\bfseries,text=gray!10!black] at (1.9,-0.42){testbed: LIBERO / CALVIN, both families};
\end{tikzpicture}
&
% ---- (b) counterfactual accuracy ----
\begin{tikzpicture}[baseline]
  \draw[->,gray] (0,0)--(3.8,0); \draw[->,gray] (0,0)--(0,2.4);
  \node[rotate=90,font=\scriptsize\bfseries,anchor=south,text=gray!10!black] at (-0.2,1.25){accuracy};
  \fill[cat-vla!55!black] (0.7,0) rectangle (1.5,0.4);
  \node[font=\scriptsize\bfseries] at (1.1,-0.22){today};
  \fill[cat-skill!38!black] (2.3,0) rectangle (3.1,1.95);
  \node[font=\scriptsize\bfseries] at (2.7,-0.22){target};
  \draw[gray,dashed,line width=0.6pt] (0,0.4)--(3.5,0.4) node[right,font=\scriptsize\bfseries,text=gray!10!black]{chance};
  \node[font=\footnotesize\bfseries] at (1.9,2.75){(b) Counterfactual accuracy};
  \node[font=\scriptsize\bfseries,text=gray!10!black] at (1.9,-0.42){testbed: WorldPrediction, LIBERO-Mem};
\end{tikzpicture}
\\[8mm]
% ---- (c) prediction-to-action fidelity gap ----
\begin{tikzpicture}[baseline]
  \draw[->,gray] (0,0)--(3.8,0); \draw[->,gray] (0,0)--(0,2.4);
  \node[rotate=90,font=\scriptsize\bfseries,anchor=south,text=gray!10!black] at (-0.2,1.1){score};
  \fill[cat-skill!52!black] (0.6,0) rectangle (1.5,2.0);
  \fill[cat-vla!52!black] (2.2,0) rectangle (3.1,0.85);
  \draw[BrickRed,{Stealth}-{Stealth},line width=0.8pt] (1.5,2.0)--(2.2,2.0);
  \node[BrickRed,font=\scriptsize\bfseries,above] at (1.85,2.0){gap};
  \draw[BrickRed,dashed,line width=0.6pt] (1.5,2.0)--(3.1,2.0);
  \node[font=\scriptsize\bfseries,align=center] at (1.05,-0.32){open-loop\\quality};
  \node[font=\scriptsize\bfseries,align=center] at (2.65,-0.32){closed-loop\\success};
  \node[font=\footnotesize\bfseries] at (1.9,2.75){(c) Prediction-to-action gap};
  \node[font=\scriptsize\bfseries,text=gray!10!black] at (1.9,-0.72){testbed: WorldSimBench, World-in-World};
\end{tikzpicture}
&
% ---- (d) per-capability contrast coverage ----
\begin{tikzpicture}[baseline]
  \draw[->,gray] (0,0)--(3.8,0); \draw[->,gray] (0,0)--(0,2.4);
  \node[rotate=90,font=\scriptsize\bfseries,anchor=south,text=gray!10!black] at (-0.2,1.25){coverage};
  \foreach [count=\i] \l/\v/\c in {{hid.}/0.3/cat-code!62!black, {c-fact}/0.1/cat-skill!48!black,
                                   {long}/0.45/ForestGreen!68!black, {soc.}/0.2/BrickRed!72!black}{
    \fill[\c] ({\i*0.78-0.2},0) rectangle ({\i*0.78+0.2},{\v});
    \node[font=\scriptsize\bfseries] at ({\i*0.78},-0.22) {\l};
  }
  \draw[BrickRed,dashed,line width=0.6pt] (0,2.0)--(3.5,2.0) node[right,font=\scriptsize\bfseries,text=BrickRed]{target};
  \node[font=\footnotesize\bfseries] at (1.7,2.75){(d) Per-capability contrast};
  \node[font=\scriptsize\bfseries,text=gray!10!black] at (1.7,-0.52){testbed: landscape audit};
\end{tikzpicture}
\end{tabular}}
\caption{\textbf{An actionable evaluation protocol for the future agenda} (Table~\ref{tab:future_metrics}).
The panels are \emph{schematic}: they show the axes and the \emph{target} shape of each measurement,
not measured results. (a)~a world-model policy should beat a matched direct-VLA baseline in
closed-loop success, and the per-capability gap is the advantage of prediction ($\delta$), on
LIBERO / CALVIN-style suites run with both families; (b)~counterfactual accuracy is near chance today
and should rise, on WorldPrediction and occlusion splits; (c)~a world model's open-loop quality far
exceeds its closed-loop task success, and that gap must be reported, on bridge suites; (d)~the
fraction of capabilities with a native VLA-vs-world-model contrast is far below target across the
landscape.}
\Description{Four schematic panels showing target measurement shapes, not real data. (a) Two curves
over capability, a world-model policy above a dashed direct-VLA baseline, with the gap shaded as the
advantage of prediction. (b) A short today bar near a chance line versus a tall target bar for
counterfactual accuracy. (c) A tall open-loop quality bar and a short closed-loop success bar with a
red gap arrow. (d) Four short per-capability contrast-coverage bars well below a dashed target line.}
\label{fig:future_protocol}
\end{figure}


\FloatBarrier
